\section*{Aktivitas Pembelajaran dan Praktikum}
\addcontentsline{toc}{section}{Aktivitas Pembelajaran dan Praktikum}

\begin{table}[H]
\centering
\caption{Rencana aktivitas Bab IV}
\label{tab:ch04-aktivitas}
\begin{tabularx}{\textwidth}{@{}llX@{}}
\toprule
\textbf{Sesi} & \textbf{Bentuk} & \textbf{Kegiatan} \\
\midrule
1 & Ceramah & Definisi formal bahasa dan tata bahasa \\
2 & Latihan & Klasifikasi tata bahasa pada hierarki Chomsky \\
3 & Diskusi & Hubungan tata bahasa, automata, dan fase kompilasi \\
4 & Lokakarya & Membuktikan bahasa tidak regular dan merancang PDA \\
\bottomrule
\end{tabularx}
\end{table}

\subsection*{Lokakarya Klasifikasi}
\begin{enumerate}[label=\arabic*.,leftmargin=*]
  \item Diberikan lima tata bahasa, tentukan tipenya dan jelaskan
        alasannya.
  \item Pilih satu bahasa regular dan satu bahasa bebas konteks.
        Tuliskan tata bahasa, automata, dan contoh string untuk
        masing-masing.
  \item Buktikan dengan lemma pompa bahwa
        \(\{a^n b^n \mid n \ge 0\}\) bukan regular.
  \item Rancang PDA untuk \(\{a^n b^{2n} \mid n \ge 0\}\).
  \item Diskusikan mengapa analisis semantik tidak cukup dilakukan
        dengan CFG.
\end{enumerate}

Bab ini tidak menambahkan increment proyek. Materi digunakan sebagai
dasar teoretis untuk Bab V--VII (leksikal) dan Bab VIII--XIII
(sintaksis).
